\begin{lemma}[Minimal deciding prefixes form a front]
If $h$ is locally constant, then $F^h$ from \noderef{DecidingFrontSet} is a
front on $\mathbb N$.
\end{lemma}
\begin{proof}
Minimality gives prefix-freeness.  Given any infinite enumeration $x$,
local constancy \noderef{LocallyConstant} supplies a deciding prefix $s$
of its range.  Repeatedly remove proper deciding prefixes; finiteness of
$s$ makes this process terminate at a member of $F^h$ still lying on the
same branch.  Thus $F^h$ is dense.  If $h$ is constant, the empty prefix is
the unique minimal member.  Otherwise, for every $n$ the branch beginning
with $n$ must meet $F^h$ in a nonempty prefix, so the base of $F^h$ is all
of $\mathbb N$.  These are exactly the three clauses of \noderef{Front}.
\end{proof}
